\documentclass[conference]{IEEEtran}
\IEEEoverridecommandlockouts

\usepackage{cite}
\usepackage{amsmath,amssymb,amsfonts}
\usepackage{algorithmic}
\usepackage{graphicx}
\usepackage{textcomp}
\usepackage{xcolor}
\usepackage{booktabs}
\usepackage{array}
\usepackage{url}
\usepackage{hyperref}

\graphicspath{{figures/}}

\def\BibTeX{{\rm B\kern-.05em{\sc i\kern-.025em b}\kern-.08em
    T\kern-.1667em\lower.7ex\hbox{E}\kern-.125emX}}

\begin{document}

\title{Uncertainty-Scaled Ensemble Transformer PPO for\\
Risk-Controlled Multi-Asset Trading under\\
Market Impact and Inferred Hidden Liquidity}

\author{\IEEEauthorblockN{Shaunak A. Rai}
\IEEEauthorblockA{\textit{Department of Computer Science and Engineering} \\
\textit{PES University}\\
Bengaluru, India \\
PES2UG23CS546}
}

\maketitle

\begin{abstract}
Deep reinforcement learning for multi-asset portfolio control routinely
ignores tail risk, sensitivity to the training seed, and execution-cost
realism in a single joint formulation, which leaves backtest gains
fragile once any one of these concerns is taken seriously. We address
this gap with a five-seed Transformer-PPO ensemble trained against a
CVaR-regularized reward, an environment augmented with a Kyle-style
linear market-impact penalty, and a volume-inferred liquidity-regime
signal built from a rolling volume z-score and an OHLC-dispersion
effective-spread proxy that acts as a public-tape surrogate for hidden
liquidity. At inference, ensemble heads are aggregated through
uncertainty-scaled position sizing $a = \mu/(1+\beta\bar\sigma)$, which
shrinks exposure precisely when seed-level disagreement is high. On a
2022--2024 out-of-sample split over five US ETFs, the liquidity-aware
ensemble (r008) reduces implementation shortfall by 14.3\%, mean
per-day market impact by 15.2\%, and turnover by 11.1\% relative to the
liquidity-naive baseline (r007). Uncertainty-scaled sizing collapses
test maximum drawdown from 58\% under naive ensemble averaging to 9\%
on the frictionless configuration --- a 6.5$\times$ reduction --- and
to 14\% on the liquidity-aware configuration, in both cases below the
21\% drawdown of passive SPY buy-and-hold. The system underperforms
passive SPY on absolute return and Sharpe; we therefore frame it as a
risk-control primitive rather than an alpha generator. Hidden liquidity
is treated throughout as an inferred public-tape proxy rather than a
tick-level dark-pool simulation. A FinBERT headline-sentiment fusion
pipeline is implemented and integrated into the state space; full
historical training is left to future work pending news-API quota
expansion.
\end{abstract}

\begin{IEEEkeywords}
Reinforcement learning, algorithmic trading, Transformer, ensemble methods,
CVaR, market impact, hidden liquidity, portfolio optimization.
\end{IEEEkeywords}

% =====================================================================
\section{Introduction}
\label{sec:intro}
Algorithmic trading systems are routinely evaluated on backtest Sharpe
ratio and ignore three properties that ultimately determine whether
they are deployable: control over tail risk, sensitivity to the random
seed used at training time, and execution cost under realistic
liquidity conditions. Reinforcement learning approaches have made
incremental progress on the first via CVaR-style risk-sensitive
objectives and on the third via friction-aware simulators that charge
spreads or impact penalties. The \emph{interaction} of these three
concerns, however, is rarely treated jointly: a method that controls
CVaR but fails silently under seed perturbation, or a method whose
backtest gains evaporate once a non-trivial impact model is switched
on, is not a deployable system. This paper takes the position that
honest evaluation of a trading-RL system requires reporting all three
together, and that doing so changes which architectural choices look
load-bearing.

We present a five-seed Transformer-PPO ensemble for multi-asset
portfolio control over five US ETFs that integrates four design
choices: (i) a CVaR-regularised reward that penalises the empirical
$5\%$-tail of recent step-rewards inside the gradient signal;
(ii) a Kyle-style market-impact penalty that charges a power-law
turnover cost scaled by realised liquidity, trained jointly with the
policy; (iii) a public-tape liquidity-regime inference signal --- a
20-day rolling volume z-score together with a smoothed
OHLC-dispersion effective-spread proxy --- that enters the state and
also drives a regime-scaled stochastic fill-noise term; and
(iv) a deployment-time uncertainty-scaled position-sizing rule
$a = \mu / (1 + \beta\bar\sigma)$ that shrinks exposure in proportion
to ensemble-head disagreement. Figure~\ref{fig:architecture} shows
where these components sit in the end-to-end pipeline.

We frame the resulting system explicitly as a \emph{risk-control
primitive} rather than as an alpha generator. On the 2022--2024
out-of-sample test split the headline configuration underperforms
passive SPY buy-and-hold on absolute return and Sharpe, but cuts
maximum drawdown by $59\%$ relative to SPY and by roughly an order of
magnitude relative to a naive ensemble-mean aggregator. Consistent
with the dark-pool honesty clause adopted throughout the paper,
``hidden liquidity'' refers to an inferred public-tape proxy for an
unobservable component of execution uncertainty, not a tick-level
dark-pool simulation; daily-bar OHLCV is incapable of supporting
stronger claims.

The contributions of this work are fivefold.
First, we release an end-to-end pipeline (Gymnasium environment,
causal Transformer encoder, five-seed PPO ensemble, uncertainty-scaled
aggregator, and impact-aware backtester) for risk-controlled
multi-asset trading, described in Sec.~\ref{sec:method}.
Second, we introduce a volume-inferred liquidity-regime signal
(rolling volume z-score and OHLC-dispersion spread proxy) used jointly
as a state feature and as the driver of a stochastic fill-noise term;
under impact-adjusted replay it cuts implementation shortfall by
$14.3\%$ relative to the frictionless baseline
(Sec.~\ref{sec:results-liquidity}).
Third, we show empirically that uncertainty-scaled position sizing is
the only configuration in our ablation that is simultaneously
positive-Sharpe, single-digit max drawdown, and robust to seed choice,
collapsing test maximum drawdown from $58\%$ under naive ensemble-mean
averaging to $9\%$ on the frictionless r007 configuration
(Sec.~\ref{sec:results-ablation}).
Fourth, we report a per-seed variance analysis showing that
validation-Sharpe ranking is a poor predictor of test rank across our
five seeds, providing direct evidence for ensemble-uncertainty methods
over single-best-seed selection in trading-RL
(Sec.~\ref{sec:results-perseed}).
Fifth, the training and evaluation protocol (five seeds $\times$
$100\text{k}$ PPO steps, val-Sharpe early stopping, frozen 2018--2024
temporal split, with full config and checkpoint snapshotting) is fully
documented for reproducibility in Sec.~\ref{sec:setup}.

% =====================================================================
\section{Related Work}
\label{sec:related}

\subsection{Reinforcement Learning for Algorithmic Trading}
Deep reinforcement learning has been applied to single-asset trading and
multi-asset portfolio control with a variety of policy classes. Deng
\textit{et al.}~\cite{deng2017drltrading} introduced a deep recurrent
direct-RL agent that learns trading signals end-to-end from raw price
series, while Jiang \textit{et al.}~\cite{jiang2017portfolio} formulated
portfolio allocation as a continuous-action MDP solved with a
deterministic policy-gradient framework over a CNN feature extractor.
Th\'eate and Ernst~\cite{theate2021tdqn} adapted DQN to discrete
long/short trading and benchmarked it against classical technical
strategies, and Zhang \textit{et al.}~\cite{zhang2020drlportfolio}
compared DQN, policy-gradient, and actor--critic agents across futures
contracts using volatility-scaled rewards. The FinRL
library~\cite{liu2020finrl} subsequently standardized environment APIs
and baselines (PPO~\cite{schulman2017ppo}, A2C, DDPG, SAC) for
reproducible research. Most of these works adopt MLP or shallow
recurrent~\cite{hochreiter1997lstm} encoders and treat execution as
frictionless or with a flat transaction-cost penalty; in contrast, our
system pairs a causal Transformer~\cite{vaswani2017attention} encoder
with a multi-asset PPO ensemble whose actions are explicitly shrunk by
inter-head disagreement, and trains under a Kyle-style impact term and
a volume-inferred liquidity regime.

\subsection{Market Microstructure and Impact Models}
Quantitative models of execution cost trace back to Kyle's
informed-trader equilibrium~\cite{kyle1985impact}, which yields the
linear price-impact coefficient $\lambda$ that we use as a training-time
penalty. Almgren and Chriss~\cite{almgren2001execution} formalized
optimal execution as a mean--variance trade-off between temporary impact
and timing risk, and Obizhaeva and Wang~\cite{obizhaeva2013impact}
extended this to dynamic, resilient limit-order-book supply. More
empirically, Cont \textit{et al.}~\cite{cont2014priceimpact} showed that
short-horizon price moves are well explained by net order-flow
imbalance, and Gatheral~\cite{gatheral2010impact} derived
no-dynamic-arbitrage constraints on impact decay kernels.
Bouchaud \textit{et al.}~\cite{bouchaud2018trades} provide the
contemporary synthesis of these results in a single textbook treatment.
Because we operate on daily bars rather than tick-level
order-book data, we deliberately use Kyle's $\lambda$ in its most
parsimonious form and treat unobserved liquidity as a public-tape
inference problem rather than a simulated dark-pool replay.

\subsection{Uncertainty Quantification in Deep RL}
Quantifying epistemic uncertainty in deep models is the basis of our
ensemble shrinkage rule. Gal and Ghahramani~\cite{gal2016dropout} cast
dropout at inference as approximate Bayesian inference, and
Lakshminarayanan \textit{et al.}~\cite{lakshminarayanan2017ensembles}
showed that simple deep ensembles are a strong, scalable alternative
that produces well-calibrated predictive variance. Within RL, Osband
\textit{et al.}~\cite{osband2016bootstrapped} used bootstrapped Q-heads
for posterior-sampling exploration, and Chua
\textit{et al.}~\cite{chua2018pets} combined probabilistic dynamics
ensembles with trajectory sampling for sample-efficient model-based
control. Prior work uses uncertainty primarily for exploration or for
model-rollout weighting. We instead apply ensemble disagreement at
deployment time as a position-sizing signal,
$a = \mu / (1 + \beta \bar\sigma)$, tying epistemic disagreement
directly to capital at risk.

\subsection{Risk-Sensitive and CVaR-Constrained RL}
Conditional Value-at-Risk (CVaR) was introduced as a coherent,
optimization-friendly tail-risk measure by Rockafellar and
Uryasev~\cite{rockafellar2000cvar}. Chow and
Ghavamzadeh~\cite{chow2014cvar} gave the first policy-gradient and
actor--critic algorithms for CVaR-constrained MDPs, and Tamar
\textit{et al.}~\cite{tamar2015cvar} derived a sampling-based CVaR
gradient that is simple to integrate into modern deep-RL pipelines.
Singh \textit{et al.}~\cite{singh2020risksensitive} more recently
combined distributional RL with risk-averse expectiles to improve
robustness under reward perturbations, and Garc\'ia and
Fern\'andez~\cite{garcia2015saferl} provide the canonical survey of
safe-RL formulations of which CVaR is one instance. We follow the
sampling-based CVaR-as-penalty formulation, with $\lambda$ tuned so
that the tail term contributes roughly an order of magnitude less
reward magnitude than mean PnL. We use the CVaR term jointly with a
Transformer encoder, microstructure-aware costs, and uncertainty-scaled
aggregation; to our knowledge this combination has not been reported
in prior CVaR-RL trading work.


% =====================================================================
\section{Problem Formulation}
\label{sec:problem}

We cast multi-asset portfolio control as a discrete-time discounted Markov
Decision Process $\mathcal{M} = (\mathcal{S}, \mathcal{A}, P, r, \gamma)$
over $N{=}5$ US ETFs at daily resolution. One environment step corresponds
to one trading day; an episode spans the full horizon of the active
split (756 days train, 252 days validation, 691 days test).

\paragraph{State space.} The observation $s_t \in \mathbb{R}^{W \times F}$
is a length-$W{=}60$ window of stacked feature vectors. The base
configuration uses $F = 65$ features per step:
$N \cdot 10 = 50$ per-asset features (OHLCV-derived: log-return,
$\log(h/c)$, $\log(l/c)$, $\log(o/c)$, 20-day volume ratio; plus five
technical indicators: RSI-14, MACD line, MACD signal, Bollinger \%B,
ATR), $\binom{N}{2} = 10$ rolling 20-day cross-asset correlation
upper-triangle entries, and $N = 5$ broadcast position weights
$a_{t-1}$. With the Phase-C liquidity-aware features enabled
(\S\ref{sec:method:impact}), two extra per-asset features --- a 20-day
rolling volume z-score and a 20-day-smoothed OHLC-dispersion spread
proxy --- raise $F$ to $75$ ($5 \times 12 + 10 + 5$). Strict causality
is enforced: $s_t$ contains no information from time $t' > t$.

\paragraph{Action space.} The action $a_t \in [-1, 1]^{N}$ is the target
portfolio weight per asset, sampled from a diagonal-Gaussian PPO actor
and clipped to the box. We deliberately omit a sum-to-one simplex
constraint: the agent may simultaneously hold long and short exposures,
and the cash position is implicit via the floor
$1 - \sum_i |a_{t,i}| \ge 0$ enforced internally. This box
parameterisation is simpler than a Dirichlet or softmax head, keeps the
policy identical to a vanilla Gaussian PPO actor, and lets the
ablations of \S\ref{sec:results-ablation} isolate the encoder and
sizing contributions without confounding from a constrained head.

\paragraph{Transition kernel.} $P(s_{t+1}\mid s_t, a_t)$ is deterministic
given the realised next-day OHLCV bar: features at $t+1$ are a
deterministic function of public market data and the new position
$a_t$. The only stochasticity in the dynamics, when Phase-C is enabled,
is the regime-scaled fill-noise term
$\varepsilon_{t,i} \sim \mathcal{N}(0, \sigma_{t,i}^{2})$ of
\eqref{eq:fillnoise} that perturbs the realised fill price.

\paragraph{Reward.} The per-step reward decomposes additively as
\begin{equation}
  r_t = \Delta\mathrm{PnL}_t \;-\; c_{\text{spread}}(a_t, a_{t-1})
        \;-\; c_{\text{impact}}(a_t, a_{t-1};\, v_t)
        \;-\; \lambda\,\widehat{\mathrm{CVaR}}_\alpha(\hat L_t),
  \label{eq:reward-mdp}
\end{equation}
where $\Delta\mathrm{PnL}_t = \sum_i a_{t-1,i}\,\log(c_{t,i}/c_{t-1,i})$
is realised log-PnL on the previously-held position, $c_{\text{spread}}$
charges a flat half-spread on turnover, $c_{\text{impact}}$ is the
Kyle's-$\lambda$ temporary-impact cost \eqref{eq:kyle}, and the final
term is a CVaR penalty \eqref{eq:cvar} computed on the rolling buffer of
recent step-losses $\hat L_t$, with $\lambda = 0.01$ and tail level
$\alpha = 0.05$. Each component is logged separately to verify that no
single term dominates the gradient (\S\ref{sec:method:sim}).

\paragraph{Discount and objective.} We use $\gamma = 0.99$, the standard
PPO~\cite{schulman2017ppo} setting at this episode length. The agent
seeks a stochastic policy $\pi_\theta(a_t \mid s_t)$ maximising
\begin{equation}
  J(\pi_\theta) = \mathbb{E}_{\pi_\theta}\!\left[\sum_{t=0}^{T-1} \gamma^t\, r_t\right],
  \quad \text{s.t.\ } a_t \in [-1,1]^{N},\ \sum_i |a_{t,i}| \le 1.
  \label{eq:objective}
\end{equation}

\begin{table}[t]
  \centering
  \caption{Notation used throughout the paper.}
  \label{tab:notation}
  \begin{tabular}{ll}
    \toprule
    Symbol & Meaning \\
    \midrule
    $N$ & Number of tradable assets ($N{=}5$) \\
    $W$ & Observation window length ($W{=}60$ days) \\
    $F$ & Features per step ($65$ base / $75$ Phase-C) \\
    $s_t \in \mathbb{R}^{W \times F}$ & State (observation window) at step $t$ \\
    $a_t \in [-1,1]^{N}$ & Target portfolio weights at step $t$ \\
    $c_{t,i}, v_{t,i}$ & Close price and traded volume of asset $i$ \\
    $\lambda, \alpha$ & CVaR weight ($0.01$) and tail level ($0.05$) \\
    $\lambda_K$ & Kyle's-$\lambda$ impact coefficient ($10^{-3}$) \\
    $\sigma_{t,i}$ & Regime-scaled fill-noise std.\ dev. \\
    $\gamma$ & Discount factor ($0.99$) \\
    $T$ & Episode horizon (756 / 252 / 691 days) \\
    $N_e$ & Ensemble size ($N_e{=}5$ seeds) \\
    \bottomrule
  \end{tabular}
\end{table}

% =====================================================================
\section{Methodology}
\label{sec:method}

This section describes the six components of the proposed system, in the
order in which an observation flows through it: the simulator and reward
(IV-A), the state representation (IV-B), the market-impact and
liquidity-aware execution model (IV-C), the Transformer encoder
(IV-D), the CVaR-regularized reward derivation (IV-E), and the
ensemble-uncertainty aggregation rule used at inference time (IV-F).
Figure~\ref{fig:architecture} summarises the end-to-end pipeline.

\begin{figure*}[!tbp]
  \centering
  \includegraphics[width=\textwidth]{architecture}
  \caption{System architecture. Data $\rightarrow$ \texttt{MultiAssetTradingEnv}
  state assembly (per-asset OHLCV + indicators + Phase-C liquidity-regime
  features + correlation matrix + positions) $\rightarrow$ Transformer encoder
  (causal, $T{=}60$, $d{=}64$) $\rightarrow$ $5\times$ PPO heads (independent
  seeds, val-Sharpe early-stop) $\rightarrow$ uncertainty-scaled aggregation
  $a = \mu / (1 + \beta \bar\sigma)$ $\rightarrow$ Kyle's-$\lambda$ market-impact
  cost and stochastic fill-noise loop back into reward.}
  \label{fig:architecture}
\end{figure*}

\subsection{Simulator Design}
\label{sec:method:sim}

The environment, \texttt{MultiAssetTradingEnv}, is a Gymnasium-compatible
discrete-time simulator over $N{=}5$ US ETFs (SPY, QQQ, IWM, GLD, TLT)
sampled at daily resolution. At each step $t$ the agent emits a
continuous action $a_t \in [-1, 1]^{5}$ interpreted as the target
portfolio weight per asset. We deliberately use an unconstrained
\texttt{Box} action space rather than a Dirichlet or softmax head: the
agent may simultaneously hold long and short exposures across assets,
and the gross portfolio weight is unconstrained at $\sum_i |a_{t,i}| \le
N$. This choice keeps the policy head identical to a vanilla Gaussian
PPO actor and avoids tying the cross-asset capital-allocation problem to
an unnecessarily restrictive simplex parameterisation.

The per-step reward decomposes additively into four terms,
\begin{equation}
  r_t = r^{\text{pnl}}_t + r^{\text{spread}}_t + r^{\text{impact}}_t
        + \lambda\,\widehat{\text{CVaR}}_\alpha(\mathbf{r}_{t-K:t}),
  \label{eq:reward}
\end{equation}
where $r^{\text{pnl}}_t = \sum_i a_{t-1,i}\,\log(c_{t,i}/c_{t-1,i})$ is
the realised log-PnL on the previously-held positions,
$r^{\text{spread}}_t = -s\,\sum_i |a_{t,i}-a_{t-1,i}|$ charges a flat
half-spread $s$ on turnover, $r^{\text{impact}}_t$ is the Kyle's-$\lambda$
temporary-impact cost defined in
Section~\ref{sec:method:impact}, and the final term is the
CVaR penalty derived in Section~\ref{sec:method:cvar}. Each component is
logged separately so that none can silently dominate the gradient signal,
and we verify post hoc that no single term contributes more than $80\%$
of $|r_t|$ in expectation.

Each per-asset feature row carries five raw OHLCV-derived signals
(log-return, $\log(h/c)$, $\log(l/c)$, $\log(o/c)$, and a 20-day
volume ratio) plus five technical indicators chosen to span momentum,
trend, and volatility regimes: a 14-period RSI (rescaled to $[-1, 1]$),
the MACD line and signal line at the standard $(12, 26, 9)$
configuration, and the Bollinger \%B and bandwidth at $(20, 2)$. The
agent also observes the rolling 20-day correlation upper triangle of
log-returns across the five assets, providing portfolio-level
co-movement awareness without requiring it to re-derive correlations
from the per-asset features.

\subsection{State Representation}
\label{sec:method:state}

The base observation is a $W \times F$ tensor with $W{=}60$ historical
steps and $F{=}65$ features per step:
\begin{itemize}
  \item \textbf{Per-asset block ($5 \times 10 = 50$):} the five OHLCV-derived
        features and five technical indicators above, concatenated for
        each of the five assets;
  \item \textbf{Cross-asset correlation block ($10$):} the upper triangle
        of the rolling 20-day log-return correlation matrix
        $\binom{5}{2} = 10$ entries;
  \item \textbf{Position block ($5$):} the agent's currently-held
        weights $a_{t-1} \in [-1,1]^{5}$, broadcast across the time
        dimension so that the encoder sees its own posture at every
        attention step.
\end{itemize}
When the Phase-C liquidity-regime features are enabled
(Sec.~\ref{sec:method:impact}), two additional features per asset are
appended, taking $F$ to $75$. All non-position features are normalised
using mean and standard deviation computed exclusively on the training
split (2018--2020); positions are already in $[-1, 1]$ and pass through
unnormalised.

\subsection{Market Impact and Liquidity-Aware Execution}
\label{sec:method:impact}

To inject the realistic frictions promised by the title we extend the
simulator with a market-impact cost and a volume-inferred
liquidity-regime signal. \emph{We frame both as public-tape
microstructure proxies for an unobservable component of execution
uncertainty rather than as a tick-level dark-pool simulator}: daily-bar
OHLCV cannot resolve hidden order flow, but it can resolve the
moments at which liquidity supply changes character, which is sufficient
to teach the agent to throttle exposure in those regimes.

\paragraph{Kyle's-$\lambda$ temporary impact.} Following the linear
price-impact tradition initiated by~\cite{kyle1985impact}, we charge a
power-law cost on each step's realised trade,
\begin{equation}
  r^{\text{impact}}_t = -\sum_i \lambda_K
       \,|a_{t,i}-a_{t-1,i}|^{\alpha}
       \,\frac{\bar{v}_i}{v_{t,i}+\varepsilon},
  \label{eq:kyle}
\end{equation}
with $\lambda_K = 10^{-3}$ and $\alpha = 1.5$. The dimensionless
liquidity factor $\bar{v}_i / v_{t,i}$ amplifies impact on
unusually thin days (where realised volume falls below the per-asset
median $\bar{v}_i$) and dampens it on heavy-volume days, giving the
agent a continuous gradient to avoid trading into illiquidity. With
these constants a full-position flip incurs roughly $10$~bps at median
liquidity and $\sim 30$~bps when volume is one third of the reference,
which matches the order of magnitude of empirical implementation
shortfalls reported for ETF block trades.

\paragraph{Liquidity-regime features.} We expose two additional per-asset
features that summarise the inferred liquidity regime. The first is a
20-day rolling volume z-score, $z_{t,i} = (v_{t,i}-\mu^v_{20,i})/
\sigma^v_{20,i}$, clipped to $[-5, 5]$: extreme negative values proxy
hidden-liquidity withdrawal (thin tape, wide spreads), and extreme
positives proxy forced unwinds or information events. The second is a
20-day-smoothed OHLC-dispersion spread proxy,
$\hat{s}_{t,i} = \mathrm{MA}_{20}\!\left[(h_{t,i}-l_{t,i})/c_{t,i}\right]$,
which is non-negative by construction and rises during illiquid regimes
when intraday ranges widen.

\paragraph{Stochastic fill noise.} To translate the inferred regime into
realised PnL dispersion, we replace the deterministic close-price fill
with a noisy fill,
\begin{equation}
  \tilde{c}_{t,i} = c_{t,i}\,(1 + \varepsilon_{t,i}), \qquad
  \varepsilon_{t,i} \sim \mathcal{N}(0, \sigma_{t,i}^{2}),
  \label{eq:fillnoise}
\end{equation}
with regime-scaled volatility
$\sigma_{t,i} = \sigma_{\text{base}}\,(1 + k\,\hat{s}_{t,i})$,
$\sigma_{\text{base}} = 5\,\text{bps}$ and $k = 10$. Realised
log-returns are then computed against $\tilde{c}_{t,i}$ rather than
$c_{t,i}$. Together (\ref{eq:kyle}) and (\ref{eq:fillnoise}) give the
agent two coupled signals --- a deterministic cost penalty and a
stochastic outcome dispersion --- both rising monotonically with the
inferred illiquidity proxies.

\subsection{Transformer Encoder}
\label{sec:method:transformer}

The features extractor is a causal Transformer~\cite{vaswani2017attention}
that maps the $W \times F$ observation window to a $d_{\text{model}}$-
dimensional latent. We use $d_{\text{model}}{=}64$, two encoder layers,
four attention heads, and a feed-forward dimensionality of $128$. A
learned positional embedding is added to the linearly-projected features,
and a strict upper-triangular causal mask prevents the encoder from
attending to future steps in the window. The pooled output of the final
step ($t = W{-}1$) is forwarded as the latent observation to the actor
and critic MLP heads. The extractor contains $\approx\!75\text{k}$
parameters and the full PPO policy $\approx\!126\text{k}$.

We adopt the pre-LayerNorm (``norm-first'') variant rather than the
original post-LN block. At this scale and with PPO's noisy advantage
estimates, post-LN training was prone to early-iteration gradient spikes
that occasionally pushed the policy into the no-trade absorbing region;
pre-LN trained reliably across all five seeds without warm-up. The
extractor is implemented as a Stable-Baselines3
\texttt{BaseFeaturesExtractor} so that the actor and critic remain the
same two-layer 128-unit MLP used by the MLP baseline; this isolates the
encoder as the only architectural change in the Transformer-vs-MLP
comparison.

\subsection{CVaR-Regularized Reward}
\label{sec:method:cvar}

To penalise tail losses inside the gradient signal rather than only at
the model-selection stage, we add a Conditional Value-at-Risk regulariser
to the reward (\ref{eq:reward}). Following the
Rockafellar--Uryasev convex
formulation~\cite{rockafellar2000cvar},
\begin{equation}
  \text{CVaR}_\alpha(L) \;=\; \min_{v\in\mathbb{R}}
       \left\{ v + \tfrac{1}{\alpha}\,
              \mathbb{E}\!\left[(L - v)^{+}\right] \right\},
  \label{eq:cvar}
\end{equation}
where $L$ denotes the loss random variable (here, $-r_t$) and
$\alpha = 0.05$ corresponds to the worst-$5\%$-tail. We follow
\cite{tamar2015cvar} in treating the inner minimisation by sampling: a
rolling buffer of the last $K{=}252$ net step-rewards is maintained, and
the empirical $\alpha$-tail mean is appended to the reward as
$\lambda\,\widehat{\text{CVaR}}_\alpha$ with $\lambda = 0.01$. Because
$\widehat{\text{CVaR}}_\alpha$ is negative whenever the agent has
realised tail losses, this term acts as a negative shaping signal on
drawdowns rather than on mean returns.

The CVaR weight $\lambda$ is itself a sensitive hyperparameter. In an
early sensitivity sweep we observed that $\lambda = 0.1$ caused the
CVaR term to account for roughly $88\%$ of the reward magnitude and
collapsed the policy onto a near-flat allocation that minimised the
penalty without trading; $\lambda = 0.01$ keeps the CVaR contribution
within the same order of magnitude as the other reward components and
recovers a learning-active policy. This calibration is reported in the
ablations of Section~\ref{sec:results}.

\subsection{Ensemble and Uncertainty-Scaled Position Sizing}
\label{sec:method:ensemble}

The final component aggregates an ensemble of $N_e = 5$ independently-
trained PPO~\cite{schulman2017ppo} heads. Each head uses a distinct
random seed drawn from $\{42, 101, 202, 303, 404\}$ controlling
parameter initialisation, environment stochasticity, and minibatch
sampling. Heads are trained with a validation-Sharpe early-stopping
callback (\texttt{ValEvalCallback}) that runs a deterministic backtest
on the held-out 2021 validation split every $10\text{k}$ environment
steps and retains the highest-Sharpe checkpoint. We find empirically
that the best-validation step varies substantially across seeds (in our
runs the per-seed best checkpoints land at $20$, $60$, $80$, $90$, and
$100\,\text{k}$ steps), so a single global early-stop schedule is not
viable; the per-seed callback is a necessary part of the protocol, not
an afterthought.

At inference time, given the per-head action distributions
$\{a^{(j)}_t\}_{j=1}^{N_e}$, we compute the per-action mean and
standard deviation across heads,
\begin{equation}
  \mu_{t,i} = \tfrac{1}{N_e}\sum_j a^{(j)}_{t,i}, \qquad
  \sigma_{t,i} = \sqrt{\tfrac{1}{N_e}\sum_j (a^{(j)}_{t,i}-\mu_{t,i})^{2}},
\end{equation}
and emit the uncertainty-scaled aggregated action
\begin{equation}
  a_t \;=\; \frac{\mu_t}{1 + \beta\,\bar{\sigma}_t},
  \qquad \bar{\sigma}_t = \tfrac{1}{N}\sum_i \sigma_{t,i},
  \label{eq:uncert}
\end{equation}
with $\beta = 5$ as the default shrinkage gain. The intuition is direct:
when the heads agree the cross-head dispersion $\bar{\sigma}_t \to 0$
and the aggregated action coincides with the simple mean; when they
disagree --- as they do persistently on out-of-distribution regimes,
where in our experiments $\bar{\sigma}_t \approx 0.77$ --- the
denominator inflates and the aggregated position shrinks toward zero.
This shrinkage rule is the principal mechanism by which the system
trades risk control for foregone alpha, and Fig.~\ref{fig:architecture}
shows where (\ref{eq:uncert}) sits in the inference pipeline.

% =====================================================================
\section{Experimental Setup}
\label{sec:setup}

\subsection{Data and Splits}
\label{sec:setup:data}
The investable universe consists of $N{=}5$ liquid US ETFs spanning equity,
size, duration, and commodity exposures: SPY (S\&P 500), QQQ (Nasdaq 100),
IWM (Russell 2000), TLT (long-duration Treasuries), and GLD (gold). All
series are daily \emph{adjusted} OHLCV bars sourced from Alpha Vantage and
cached locally as Parquet under \texttt{data/}; using adjusted closes
ensures that splits and dividends do not appear as spurious one-day
returns. We adopt a strict temporal three-way split with no shuffling:
2018--2020 for training (756 trading days), calendar 2021 for validation
(252 days), and 2022--2024 for the held-out test set (691 days). Per
CLAUDE.md \S6.2, normalisation statistics (per-feature mean and standard
deviation) are fitted exclusively on the training window and applied
unchanged to validation and test, and a unit test asserts that the state
constructed at time $t$ contains no field derived from $t{+}1$ or later.

\subsection{Baselines}
\label{sec:setup:baselines}
We compare the proposed system against six external baselines and five
internal RL ablation baselines, all evaluated on the same 2022--2024 OOS
split. The external baselines are:
(i) \emph{SPY buy-and-hold}, a single-asset passive long;
(ii) \emph{equal-weight 5-ETF}, a daily-rebalanced uniform allocation
across the universe;
(iii) \emph{60/40 SPY/TLT}, a daily-rebalanced classical balanced
portfolio;
(iv) \emph{MA-crossover}, a 50/200-day SMA crossover on SPY with a
one-day-lagged signal to avoid look-ahead;
(v) \emph{TWAP}, a 10-day linear ramp toward the equal-weight target,
modelling time-sliced execution; and
(vi) \emph{VWAP-proxy}, equal-weight target executed at the daily OHLC4
midprice rather than the close, modelling volume-weighted execution.
The internal RL baselines are MLP+CVaR (run \texttt{r005}), the
single-seed Transformer at the early-stopping checkpoint
(\texttt{r006-ckpt50k}), the naive-mean five-seed ensemble
(\texttt{r007 mean}), the uncertainty-scaled ensemble \texttt{r007}
(frictionless training), and the uncertainty-scaled ensemble
\texttt{r008} (Phase-C liquidity-aware training).

\subsection{Metrics}
\label{sec:setup:metrics}
We report six standard return/risk metrics: annualised return,
Sharpe ratio (with $r_f{=}0$), Sortino ratio (downside-only denominator),
maximum drawdown, Calmar ratio (annualised return divided by max
drawdown), and Conditional Value-at-Risk at the $5\%$ tail of daily
returns. All are computed by QuantStats on daily NAV. For the Phase-C
execution comparison (Sec.~\ref{sec:results-liquidity}) we additionally
report three implementation-shortfall metrics: total implementation
shortfall in basis points, mean per-day market impact in basis points
(from Eq.~\ref{eq:kyle}), and mean daily portfolio turnover.

\subsection{Hyperparameters}
\label{sec:setup:hyper}

\begin{table}[t]
  \centering
  \caption{Hyperparameter settings.}
  \label{tab:hyper}
  \footnotesize
  \begin{tabular}{ll@{\hskip 1.5em}ll}
    \toprule
    \multicolumn{2}{l}{\textit{PPO}} & \multicolumn{2}{l}{\textit{Transformer}} \\
    lr                & $3{\times}10^{-4}$              & $d_{\text{model}}$   & $64$ \\
    $n_{\text{steps}}$& $256$                           & layers               & $2$ \\
    minibatch         & $64$                            & heads                & $4$ \\
    epochs            & $10$                            & FFN dim              & $128$ \\
    $\gamma$          & $0.99$                          & dropout              & $0.1$ \\
    GAE $\lambda$     & $0.95$                          & window $T$           & $60$ \\
    clip              & $0.2$                           & norm                 & pre-LN \\
    $c_{\text{ent}}$  & $0.01$                          & pos.\ embed          & learned \\
    $c_v$             & $0.5$                           & head                 & MLP $[128,128]$ \\
    \midrule
    \multicolumn{2}{l}{\textit{Ensemble}} & \multicolumn{2}{l}{\textit{CVaR \& Phase-C}} \\
    seeds             & 5 (42--404)                     & CVaR $\lambda$       & $0.01$ \\
    steps             & $100$k                          & CVaR $\alpha$        & $0.05$ \\
    eval interval     & $10$k                           & buffer $K$           & $252$ \\
    select            & best val Sharpe                 & Kyle $\lambda_K$     & $10^{-3}$ \\
    $\beta$           & $5$                             & Kyle exp.\ $\alpha$  & $1.5$ \\
    half-spread $s$   & $1$ bp                          & $\sigma_{\text{base}}$, $k$ & $5$ bps, $10$ \\
    \bottomrule
  \end{tabular}
\end{table}

Table~\ref{tab:hyper} consolidates all hyperparameters used in the
reported runs. PPO settings follow the Stable-Baselines3 defaults except
for the rollout length $n_{\text{steps}}{=}256$, which trades a smaller
on-policy buffer for more frequent policy updates at our $100\text{k}$-step
budget; we therefore deviate from the SB3 default of $2048$. The
entropy coefficient is set to $0.01$ rather than $0$ to maintain
exploration on the continuous portfolio-weight action space. The
Transformer extractor contains $\approx\!75\text{k}$ parameters and the
full PPO policy $\approx\!126\text{k}$. Every seed controls numpy,
torch, Python \texttt{random}, and the Gym environment RNG
simultaneously, and the entire seed set $\{42, 101, 202, 303, 404\}$
is logged in each run's config snapshot.

\subsection{Hardware and Reproducibility}
\label{sec:setup:hw}
Smoke tests, single-seed development runs, all six external baselines,
and impact-adjusted backtest replays were executed locally on an AMD
Ryzen 9 9950X3D + NVIDIA RTX 5080 workstation. Five-seed ensemble
training (runs \texttt{r007} and \texttt{r008}) was offloaded to a
Google Colab NVIDIA T4 (16 GB) instance, consuming approximately
$5\text{--}6$ GPU-hours per five-seed run at $100\text{k}$ steps each.
Per CLAUDE.md \S6.4, every run writes a timestamped directory under
\texttt{runs/} containing the full config snapshot, the seed, the
\texttt{ValEvalCallback} log, all $10\text{k}$-step checkpoints, and the
final TensorBoard event file, so that each row in
Tables~\ref{tab:main}--\ref{tab:liquidity} traces back to an exact
reproducible artefact. Point-estimate Sharpe ratios in
Table~\ref{tab:main} are accompanied by $95\%$ bootstrap confidence
intervals computed via circular block bootstrap ($N{=}2000$,
block${=}21$ trading days) to account for return autocorrelation.

% =====================================================================
\section{Results}
\label{sec:results}

\begin{figure}[t]
  \centering
  \includegraphics[width=\columnwidth]{fig01_equity_curves.pdf}
  \caption{Equity curves: r007/r008 ensembles vs.\ baselines on 2022--2024 OOS.}
  \label{fig:equity}
\end{figure}

\begin{figure}[t]
  \centering
  \includegraphics[width=\columnwidth]{fig05_ensemble_comparison.pdf}
  \caption{Ensemble mean vs.\ uncertainty-scaled aggregation: the
  shrinkage on disagreement collapses test max drawdown from
  $0.58$ to $0.09$ on r007.}
  \label{fig:uncertainty}
\end{figure}

\begin{figure}[t]
  \centering
  \includegraphics[width=\columnwidth]{fig09_liquidity_stress.pdf}
  \caption{Phase-C contribution: r007 vs.\ r008 implementation shortfall,
  mean impact, mean spread, mean turnover (impact-adjusted replay).}
  \label{fig:liquidity}
\end{figure}

We report all results on the held-out 2022--2024 test split (691 trading
days), using a five-seed ensemble (seeds $\{42, 101, 202, 303, 404\}$)
trained with val-Sharpe early stopping on calendar year 2021. Two
configurations of the full pipeline are compared throughout: \emph{r007}
(frictionless training) and \emph{r008} (Phase-C liquidity-aware training
with volume-z-score and OHLC-dispersion spread-proxy features plus
stochastic fill noise).

\subsection{Main Results vs.\ Benchmarks}
\label{sec:results-main}

Table~\ref{tab:main} compares the r007 uncertainty-scaled ensemble against
three passive baselines on the 691-day OOS split. The strategy
\emph{underperforms} SPY buy-and-hold on absolute return (1.27\% vs.\
11.14\%) and on Sharpe (0.27 vs.\ 0.70), and we make no claim of
out-performance against passive equities. The defensible result is risk
control: maximum drawdown is 8.64\% versus 21.28\% for SPY
($-59\%$), 21.31\% for the equal-weight 5-ETF basket ($-59\%$), and
21.99\% for a daily-rebalanced 60/40 SPY/TLT mix ($-61\%$). Calmar
(annualized return over max drawdown) at 0.147 also exceeds the 60/40
mix's 0.112 despite a near-identical Sharpe, since the strategy buys its
small return at roughly one third the drawdown. Equity trajectories are
shown in Fig.~\ref{fig:equity}. Bootstrap $95\%$ CIs (Table~\ref{tab:main}, inline) show no strategy
achieves statistically significant Sharpe separation over 691 days,
consistent with post-2022 regime difficulty. The ensemble's contribution
is better characterised by its maximum drawdown ($8.64\%$), an order of
magnitude below the naive ensemble mean ($44.2\%$, row~E of
Table~\ref{tab:ablation}) and robust across bootstrap samples.

\begin{table}[t]
  \centering
  \caption{Main results on the 2022--2024 OOS test split (691 days).}
  \label{tab:main}
  \begin{tabular}{lrrrr}
    \toprule
    Strategy & Ann.\ Ret.\ & Sharpe & Max DD & Calmar \\
    \midrule
    r007 ens.\ (ours)        & $+1.27\%$  & $+0.27\ [-1.05,{+}1.59]$ & $\mathbf{8.64\%}$ & $+0.147$ \\
    SPY buy-and-hold         & $+11.14\%$ & $+0.70\ [-0.44,{+}2.16]$ & $21.28\%$         & $+0.524$ \\
    Equal-weight 5-ETF       & $+6.36\%$  & $+0.51\ [-0.70,{+}1.89]$ & $21.31\%$         & $+0.298$ \\
    60/40 SPY/TLT            & $+2.47\%$  & $+0.25\ [-1.06,{+}1.70]$ & $21.99\%$         & $+0.112$ \\
    \bottomrule
  \end{tabular}
\end{table}

\subsection{Ablation Study}
\label{sec:results-ablation}

Table~\ref{tab:ablation} reports the full component ablation, rows A--G,
each evaluated on the same 2022--2024 OOS split. Rows A and B are
single-seed MLP runs (731 days, includes 2021 contiguous test segments
from the multi-asset pre-CVaR sweep); rows C--G are 691-day
post-2022 test runs. The drawdown trajectory is monotone-on-improvement
from row C onward: $50.9\% \to 51.4\% \to 71.8\% \to 56.3\% \to 44.2\%
\to \mathbf{8.6\%} \to 13.7\%$. Adding CVaR (B) holds drawdown roughly
constant while nudging Sharpe up; introducing the Transformer without
early stopping (C) overfits and inflates drawdown to 71.8\%; val-Sharpe
early stopping (D) recovers a strong single-seed Sharpe of 0.74 but at
56.3\% drawdown that is entirely seed-luck (see
Section~\ref{sec:results-perseed}); naive ensemble averaging (E)
\emph{worsens} Sharpe to $-0.20$ as opposing positions cancel and spread
is paid on the residual; and \emph{uncertainty-scaled sizing} (row F)
collapses drawdown by an order of magnitude to 8.6\% while flipping
Sharpe positive. Row F is the headline configuration of the paper.

Row G adds the Phase-C liquidity-aware environment (volume-regime state
features and stochastic fill noise) to the same uncertainty-sizing
ensemble. Test Sharpe regresses to $-0.27$ and test max drawdown widens
to 13.7\%, but the underlying microstructure mechanism is validated
under impact-adjusted replay (Section~\ref{sec:results-liquidity}).

\begin{table}[t]
  \centering
  \caption{Component ablation on the OOS test split. Rows A--B
  use 731 days; rows C--G use 691 days (post-2022).}
  \label{tab:ablation}
  \begin{tabular}{clrrrr}
    \toprule
    Row & Configuration & Ann.\ Ret. & Sharpe & Max DD & Calmar \\
    \midrule
    A & MLP, no CVaR             & $-3.2\%$  & $+0.17$  & $50.9\%$ & $-0.06$ \\
    B & + CVaR ($\lambda{=}0.01$)& $-1.3\%$  & $+0.18$  & $51.4\%$ & $-0.03$ \\
    C & + Transformer, no ES     & $-10.8\%$ & $+0.09$  & $71.8\%$ & $-0.15$ \\
    D & + early stop (50k)       & $+29.0\%$ & $+0.74$  & $56.3\%$ & $+0.52$ \\
    E & + 5-seed ens.\ (mean)    & $-8.8\%$  & $-0.20$  & $44.2\%$ & $-0.20$ \\
    \textbf{F} & \textbf{+ uncert.\ sizing (FULL, r007)} & $\mathbf{+1.3\%}$  & $\mathbf{+0.27}$  & $\mathbf{8.6\%}$ & $\mathbf{+0.15}$ \\
    G & + liquidity env (r008)   & $-2.3\%$  & $-0.23$  & $13.7\%$ & $-0.18$ \\
    \bottomrule
  \end{tabular}
\end{table}

\subsection{Per-Seed Variance and the Necessity of Uncertainty Sizing}
\label{sec:results-perseed}

The deceptively simple ensemble configuration hides a large per-seed
test-time variance that motivates the entire uncertainty-sizing
contribution. Table~\ref{tab:perseed} reports each r007 head evaluated
in isolation. Test Sharpe spans $[-1.00, +0.74]$ and annualized return
spans $[-54.7\%, +26.0\%]$ across the five seeds. Critically,
\emph{validation Sharpe is a poor predictor of test rank}: the best-val
seed (303, val Sharpe $1.63$) finished third on test, while the
best-test seed (101, test Sharpe $+0.74$) had only the fourth-best val
Sharpe. Picking a single head by val rank is close to a coin flip on
this dataset.

\begin{table}[t]
  \centering
  \caption{Per-seed r007 test performance (each head evaluated alone).}
  \label{tab:perseed}
  \begin{tabular}{crrrr}
    \toprule
    Seed & Val Sharpe & Test Sharpe & Test Ann.\ Ret. & Test Max DD \\
    \midrule
    42  & $+1.37$ & $-0.35$         & $-30.4\%$         & $71.1\%$ \\
    101 & $+0.68$ & $\mathbf{+0.74}$ & $\mathbf{+26.0\%}$ & $42.0\%$ \\
    202 & $+0.91$ & $-1.00$         & $-54.7\%$         & $90.4\%$ \\
    303 & $+1.63$ & $+0.02$         & $-15.2\%$         & $66.2\%$ \\
    404 & $+0.59$ & $+0.09$         & $-2.4\%$          & $31.3\%$ \\
    \midrule
    mean of heads & --- & $-0.10$ & $-15.3\%$ & $\approx 60\%$ \\
    \bottomrule
  \end{tabular}
\end{table}

The naive ensemble mean (row E of Table~\ref{tab:ablation}) is in fact
\emph{worse} than the average individual Sharpe ($-0.20$ vs.\ $-0.10$),
because heads take opposing long/short positions that cancel and the
agent pays spread on the residual. By contrast, uncertainty-scaled
sizing -- which shrinks the position by $1/(1+\beta\bar\sigma)$ when
heads disagree -- drives mean head standard deviation $\bar\sigma
\approx 0.77$ on test, indicating that disagreement is sustained across
the unseen 2022 bear, 2023 rally, and 2024 mixed regimes. This
shrinkage is the mechanism that takes the ensemble from $0.58$ to
$0.09$ in test max drawdown (Fig.~\ref{fig:uncertainty}). Row F is the
only configuration that is simultaneously (a) positive Sharpe, (b)
single-digit max drawdown, and (c) not contingent on a lucky seed pick;
we therefore present it as the only \emph{deployable} variant.

\subsection{Liquidity-Aware Execution Effects (Phase C, r008)}
\label{sec:results-liquidity}

To evaluate whether the Phase-C liquidity-regime state features and
stochastic fill noise actually shape execution behaviour, we replay the
r007 and r008 ensemble-uncertainty policies through the same
impact-adjusted environment (Kyle's-$\lambda$ enabled) on the
2022--2024 test split. Table~\ref{tab:liquidity} reports the four
execution-quality metrics. All four improve in the predicted direction:
implementation shortfall falls from 719.1 bps to 616.6 bps
($-14.3\%$, roughly 37 bps/yr saved over the 691-day window); mean
daily impact falls from 0.79 bps to 0.67 bps ($-15.2\%$); mean daily
spread falls from 0.25 bps to 0.22 bps ($-12.0\%$); and mean daily
turnover falls from 0.2478 to 0.2203 ($-11.1\%$). The agent has
learned to trade less in inferred-illiquid regimes, which is precisely
the behaviour the Phase-C training signal was designed to induce
(Fig.~\ref{fig:liquidity}).

\begin{table}[t]
  \centering
  \caption{Impact-adjusted replay: r007 vs.\ r008 execution metrics on
  the 2022--2024 test split, ensemble-uncertainty policy in both rows.}
  \label{tab:liquidity}
  \begin{tabular}{lrrr}
    \toprule
    Metric                          & r007 & r008 & $\Delta$ \\
    \midrule
    Impl.\ shortfall (bps)          & $719.1$  & $\mathbf{616.6}$  & $\mathbf{-14.3\%}$ \\
    Mean daily impact (bps)         & $0.79$   & $\mathbf{0.67}$   & $\mathbf{-15.2\%}$ \\
    Mean daily spread (bps)         & $0.25$   & $\mathbf{0.22}$   & $\mathbf{-12.0\%}$ \\
    Mean daily turnover             & $0.2478$ & $\mathbf{0.2203}$ & $\mathbf{-11.1\%}$ \\
    \bottomrule
  \end{tabular}
\end{table}

The cost of this microstructure improvement is that gross Sharpe
contracted from $+0.27$ to $-0.23$ between r007 and r008. We attribute
this to insufficient training budget: at 100k PPO steps per seed, the
liquidity-aware policy reduced gross turnover but had not yet converted
that reduction into higher per-trade quality. A net-of-impact
comparison still favours r007 on this budget ($\approx -73$ bps/yr vs.\
$\approx -400$ bps/yr), but the Phase-C \emph{mechanism} -- lower
turnover, lower impact, lower implementation shortfall -- is
unambiguously validated by Table~\ref{tab:liquidity}. We discuss the
training-budget gap as a limitation in Section~\ref{sec:discussion}.

% =====================================================================
\section{Discussion}
\label{sec:discussion}

\subsection{What Worked, What Didn't}

Three results held up. Row~F of Table~\ref{tab:ablation} is the only
configuration with single-digit test max drawdown ($8.6\%$); every
other variant, including the early-stopped single seed (row~D,
$56.3\%$) and naive ensemble mean (row~E, $44.2\%$), sits in
$44$--$72\%$. The aggregator $a = \mu/(1+\beta\bar\sigma)$ is best
read as a risk-control primitive that buys bounded drawdown by paying
foregone return when heads disagree OOD. Under impact-adjusted replay
(Table~\ref{tab:liquidity}) r008 cuts implementation shortfall by
$14.3\%$, impact by $15.2\%$, and turnover by $11.1\%$, confirming
that the liquidity-regime features functioned as intended. And the
per-seed best-validation steps spanned $\{20, 60, 80, 90, 100\}$k:
per-seed checkpointing is load-bearing, not incidental.

Two results did not. The headline configuration earns $+1.27\%$
annualized at Sharpe $+0.27$ against SPY's $+11.14\%$/$+0.70$
(Table~\ref{tab:main}); the only honest comparative claim is the
$59\%$ drawdown reduction. Adding the Phase-C environment dropped
r008 test Sharpe to $-0.27$ at $13.7\%$ drawdown (rows~F, G); we
attribute this to a training-budget gap in which $100$k PPO steps
shrank turnover but had not yet converted that into per-trade quality
-- a budget finding, not a refutation of liquidity awareness, since
the mechanism is validated by Table~\ref{tab:liquidity}.

\subsection{Limitations}
\label{sec:disc:limits}

The simulator operates on daily OHLCV bars without a tick-level order
book; per Sec.~\ref{sec:method:impact}, hidden liquidity is treated as
an unobservable execution-uncertainty component inferred from public-tape
proxies (volume z-score, dispersion spread), not a simulated dark-pool
replay. Validation uses 2021 only, so the row-F val Sharpe of $1.98$
carries a wide confidence interval; we anchor all claims on the OOS
test Sharpe ($+0.27$). Frictions are simplified: a flat $1$~bps
half-spread plus the Phase-C Kyle's-$\lambda$ impact and regime-scaled
fill noise, with no borrow cost, locate fees, or hard-to-borrow
exclusions. All results are backtested; no live or paper-brokerage
deployment was attempted. Finally, the r009 FinBERT sentiment pipeline
(fetch, inference, env integration, unit tests) is implemented but
untrained: free-tier news-API quotas blocked the multi-year backfill
needed for end-to-end training; we report this as future work.

% =====================================================================
\section{Conclusion and Future Work}
\label{sec:conclusion}
We presented a 5-seed Transformer-PPO ensemble trained with
CVaR$_{5\%}$ regularization, a Kyle's-$\lambda$ impact penalty, and
volume-inferred liquidity-regime features, aggregated via
$a = \mu / (1 + \beta\,\bar\sigma)$. The aggregator is the
load-bearing component: it collapses test max drawdown from $58\%$
(naive mean) to $9\%$ on r007 and from $60\%$ per-seed to $14\%$ on
r008, and is the only configuration that survives seed variance. We
frame the contribution as a \emph{drawdown-control primitive}, not a
passive-equity-beating alpha system. Future work: completing r009
FinBERT sentiment fusion (blocked on news-API backfill), replacing
daily-bar Kyle's-$\lambda$ with intra-day LOB replay or
Almgren--Chriss execution, a \texttt{MarketConfig} abstraction for
Indian equities (NSE/BSE), and Interactive Brokers paper-trading
deployment.

% =====================================================================
{\footnotesize
\bibliographystyle{IEEEtran}
\bibliography{references}
}

\end{document}
